\documentclass[11pt]{article}
\usepackage[utf8]{inputenc}
\usepackage{amsmath,amssymb}
\usepackage{hyperref}
\title{Robust Climate Model Downscaling Deep Learning under Distribution Shift}
\author{Anonymous}
\date{}
\begin{document}
\maketitle

\begin{abstract}
This paper studies Climate Model Downscaling Deep Learning. Grounded in a real, citable literature \cite{ref1,ref2,ref3,ref4,ref5,ref6,ref7,ref8},
we propose a focused method and report \textbf{illustrative} experiments.
All references below are real and verifiable (OpenAlex/arXiv); the reported
numbers are simulated to illustrate intended behaviour.
\end{abstract}

\section{Introduction}
Climate Model Downscaling Deep Learning is an active research area. This work targets a concrete gap identified from
the recent literature and contributes a method together with an evaluation protocol.

\section{Related Work (real literature)}
The following works, retrieved from public bibliographic APIs, frame our study:
\begin{itemize}
\item Tripathy et al. (2023) — "Deep learning in hydrology and water resources disciplines: concepts, methods, applications, and research directions", Journal of Hydrology (cited 300).
\item Baño‐Medina et al. (2020) — "Configuration and intercomparison of deep learning neural models for statistical downscaling", Geoscientific model development (cited 260).
\item Wang et al. (2021) — "Deep Learning for Daily Precipitation and Temperature Downscaling", Water Resources Research (cited 222).
\item Li et al. (2019) — "Spatiotemporal imputation of MAIAC AOD using deep learning with downscaling", Remote Sensing of Environment (cited 123).
\item Sun et al. (2020) — "Statistical downscaling of daily temperature and precipitation over China using deep learning neural models: Localization and comparison with other methods", International Journal of Climatology (cited 113).
\item Rajesh et al. (2022) — "Impact of climate change on river water temperature and dissolved oxygen: Indian riverine thermal regimes", Scientific Reports (cited 110).
\end{itemize}

\section{Method}
We model distribution shift explicitly and optimise a worst-case objective over an uncertainty set of plausible test distributions. We instantiate this idea for Climate Model Downscaling Deep Learning and analyse its cost/quality trade-off.

\section{Experiments (Illustrative)}
\begin{center}
\begin{tabular}{lcc}
System & Primary Metric & Efficiency \\ \hline
Strong baseline & 0.812 & 1.00$\times$ \\
Ablation & 0.828 & 0.92$\times$ \\
\textbf{Ours} & \textbf{0.851} & \textbf{0.71$\times$}
\end{tabular}
\end{center}
\emph{Metrics simulated to illustrate the intended effect; not measured.}

\section{Conclusion}
We connected a real literature to a concrete method for Climate Model Downscaling Deep Learning. Future work will
validate the illustrative results with real experiments.

\begin{thebibliography}{9}
\bibitem{ref1} Kumar Puran Tripathy, Ashok K. Mishra. Deep learning in hydrology and water resources disciplines: concepts, methods, applications, and research directions. \emph{Journal of Hydrology}, 2023. DOI:10.1016/j.jhydrol.2023.130458
\bibitem{ref2} Jorge Baño‐Medina, Rodrigo Manzanas, José Manuel Gutiérrez. Configuration and intercomparison of deep learning neural models for statistical downscaling. \emph{Geoscientific model development}, 2020. DOI:10.5194/gmd-13-2109-2020
\bibitem{ref3} Fang Wang, Di Tian, Lisa L. Lowe et al.. Deep Learning for Daily Precipitation and Temperature Downscaling. \emph{Water Resources Research}, 2021. DOI:10.1029/2020wr029308
\bibitem{ref4} Lianfa Li, Meredith Franklin, Mariam Girguis et al.. Spatiotemporal imputation of MAIAC AOD using deep learning with downscaling. \emph{Remote Sensing of Environment}, 2019. DOI:10.1016/j.rse.2019.111584
\bibitem{ref5} Lei Sun, Yufeng Lan. Statistical downscaling of daily temperature and precipitation over China using deep learning neural models: Localization and comparison with other methods. \emph{International Journal of Climatology}, 2020. DOI:10.1002/joc.6769
\bibitem{ref6} M. Rajesh, S. Rehana. Impact of climate change on river water temperature and dissolved oxygen: Indian riverine thermal regimes. \emph{Scientific Reports}, 2022. DOI:10.1038/s41598-022-12996-7
\bibitem{ref7} Neelesh Rampal, Peter B. Gibson, Abha Sood et al.. High-resolution downscaling with interpretable deep learning: Rainfall extremes over New Zealand. \emph{Weather and Climate Extremes}, 2022. DOI:10.1016/j.wace.2022.100525
\bibitem{ref8} Yongjian Sun, Kefeng Deng, Kaijun Ren et al.. Deep learning in statistical downscaling for deriving high spatial resolution gridded meteorological data: A systematic review. \emph{ISPRS Journal of Photogrammetry and Remote Sensing}, 2024. DOI:10.1016/j.isprsjprs.2023.12.011
\end{thebibliography}
\end{document}
